\documentclass[../../main]{subfiles}

\begin{document}

\section{Corpos}
\label{sec:matematica:algebra:estruturas-algebricas:corpos}

Nesta seção consideramos estruturas algébricas sobre a assinatura
\[
  \Sigma = \{+, \cdot : 2\},
\]
isto é, as mesmas operações de um anel comutativo unitário, acrescidas da exigência de inversibilidade multiplicativa para elementos não nulos.

\subsection{Definição de corpo}

Um corpo é um anel comutativo unitário $(F,+,\cdot)$ tal que:

\begin{itemize}
  \item $(F,+,\cdot)$ satisfaz todos os axiomas de anel comutativo unitário;

  \item existe um elemento $1 \in F$, com $1 \neq 0$, tal que:
    \[
      \forall a \in F,\quad a \cdot 1 = 1 \cdot a = a;
    \]

  \item todo elemento não nulo possui inverso multiplicativo:
    \[
      \forall a \in F \setminus \{0\},\ \exists a^{-1} \in F \text{ tal que }
      a \cdot a^{-1} = a^{-1} \cdot a = 1.
    \]
\end{itemize}

\subsection{Estrutura multiplicativa}

Em um corpo $(F,+,\cdot)$, o conjunto $F \setminus \{0\}$ munido da multiplicação forma um grupo abeliano:
\[
  (F \setminus \{0\}, \cdot).
\]

Isto significa que:

\begin{itemize}
  \item a multiplicação é associativa;
  \item existe elemento neutro $1$;
  \item todo elemento não nulo é invertível;
  \item a multiplicação é comutativa.
\end{itemize}

\subsection{Estrutura aditiva}

O par $(F,+)$ é um grupo abeliano, isto é:

\begin{itemize}
  \item a adição é associativa e comutativa;
  \item existe elemento neutro $0$;
  \item todo elemento possui inverso aditivo.
\end{itemize}

\subsection{Domínio de integridade}

Um anel comutativo unitário $(A,+,\cdot)$ é um domínio de integridade se:
\[
  \forall a,b \in A,\quad a \cdot b = 0 \Rightarrow a = 0 \ \text{ou}\ b = 0.
\]

\subsection{Relação entre corpos e domínios de integridade}

Todo corpo é um domínio de integridade.

\textit{Demonstração:}
Se $a \cdot b = 0$ e $a \neq 0$, então existe $a^{-1}$ e:
\[
  a^{-1}(a \cdot b) = (a^{-1}a)\cdot b = 1 \cdot b = b,
\]
logo $b = 0$.

\subsection{Exemplos fundamentais}

\begin{itemize}
  \item $(\mathbb{Q}, +, \cdot)$ é um corpo;
  \item $(\mathbb{R}, +, \cdot)$ é um corpo;
  \item $(\mathbb{C}, +, \cdot)$ é um corpo;
  \item $(\mathbb{Z}_p, +, \cdot)$ é um corpo se $p$ é primo.
\end{itemize}

\subsection{Observação estrutural}

Um corpo pode ser visto como um anel comutativo unitário no qual a operação multiplicativa restrita a $F \setminus \{0\}$ forma um grupo abeliano.

Essa propriedade torna possível a noção de divisão entre elementos não nulos e é a base para a construção de espaços vetoriais.

\end{document}
